\chapter{Bases Teóricas e Tecnológicas}
\label{cap:fundamentos}

Este capítulo deve descrever a base conceitual do TCC. Explicar os pilares teóricos específicos que sustentam o projeto, como por exemplo: arquitetura de software, requisitos, ciclo de vida, etc. 

Apresente uma visão geral do tópico trabalhado no TCC, sem tentar falar de Engenharia de Software como um todo. Pode falar do ciclo de vida de processos de software, mas o  foco deve ser somente no que foi \textbf{efetivamente trabalhado e realizado no seu próprio TCC}.

O capítulo também deve fazer uma síntese dos principais conceitos instrumentais (ferramentas e tecnologias) aplicados ao longo do TCC, destacando como eles ajudam a compreensão da experiência relatada neste documento. 

Este capítulo deve estabelecer uma conexão clara entre a teoria utilizada, as ferramentas aplicadas e as atividades práticas desenvolvidas, preparando o leitor para o capítulo seguinte, que é o próprio relato da experiência.

\section{Tópico Específico trabalhado no TCC}

Apresente os conceitos fundamentais do tópico específico da Engenharia de Software que foi trabalhado no TCC, destacando sua importância para projetos de Engenharia de Software que precisam utilizar métodos, práticas e técnicas para o desenvolvimento e a evolução de sistemas computacionais. 

Coloque no título desta seção o tópico específico da Engenharia de Software que foi o assunto principal trabalhado no TCC. Por exemplo: Teste de Software; Requisitos de Software; Integração de Software; Técnicas de Validação de Software; Implementação de Interfaces Humano-Computador; Implementação de Funções de \textit{Backend} de Software. Se trabalhou em mais de um tópico, crie uma subseção para cada tópico, e não misture dois ou mais assuntos diferentes na mesma seção.

Discuta os objetivos do tópico trabalhado, seus princípios e os desafios inerentes à criação e manutenção de software de qualidade. A seção deve estabelecer a base conceitual necessária para compreender qualquer atividade realizada no contexto do projeto descrito no trabalho.

\section{Processos de Ciclo de Vida de Software}

Aqui são apresentados os principais processos de Engenharia de Software relacionados ao tópico trabalhado. Por exemplo, processos de Verificação e Validação, se o tópico é Teste de Software; ou processos de Planejamento e Monitoramento de projeto, se o tópico é Gerência de Projeto de Software. 

São discutidas as características de cada processo, suas vantagens e limitações, bem como sua aplicabilidade em diferentes contextos. Esta fundamentação permite situar a atividade realizada dentro de um processo mais amplo de ciclo de vidade de software, que contempla o desenvolvimento e a manutenção de software.

Cada processo deve ser apresentado de forma conceitual, indicando suas atividades genéricas e permitindo relacionar a teoria de Engenharia de Software com as tarefas específicas desempenhadas durante a experiência prática, independentemente da natureza dessas tarefas.

\section{Organização do Trabalho Colaborativo em Software}

Discuta os aspectos relacionados ao trabalho em equipe no contexto de projetos de software, considerando especificamente o papel ou função desempenhado durante o trabalho de TCC. 

Aborde temas como comunicação, coordenação, divisão de responsabilidades, interfaces com outros papéis típicos em equipes técnicas e desafios comuns enfrentados por equipes de desenvolvimento relacionadas ao papel desempenhado. 

Esta seção deve fundamentar a sua análise da experiência ao participar de uma equipe real, destacando fatores humanos e organizacionais que influenciaram sua atuação no projeto.

\section{Ferramentas e Tecnologias Adotadas}

Apresente e analise as ferramentas utilizadas para apoiar os processos e atividades discutidas anteriormente. 

As ferramentas devem ser comparadas a outras amplamente utilizadas na indústria de software, justificando sua escolha para o trabalho realizado. 

Evite fazer apenas uma propaganda das ferramentas. Discuta os pontos fortes, mas analise também as limitações de cada ferramenta.
Devem ser discutidos artefatos gerados nas atividades dos processos e a forma como as ferramentas apoiam a criação e evolução destes artefatos.

A discussão de ferramentas permite contextualizar a forma como o trabalho da equipe e do autor do TCC foi organizado e como as atividades do autor se inseriram nos processos colaborativos da Engenharia de Software, conforme discutido na seção anterior.
